%%%%%%%%%%%%%%%%%%%%%%%%%%%%%%%%%%%%%%%%%%%%%%%%%%%%%%%%%%%%%%%%%%%%%%%%

%%% LaTeX Template for ECAI Papers 
%%% Prepared by Ulle Endriss (version 1.0 of 2023-12-10)

%%% To be used with the ECAI class file ecai.cls.
%%% You also will need a bibliography file (such as mybibfile.bib).

%%%%%%%%%%%%%%%%%%%%%%%%%%%%%%%%%%%%%%%%%%%%%%%%%%%%%%%%%%%%%%%%%%%%%%%%

%%% Start your document with the \documentclass{} command.
%%% Use the first variant for the camera-ready paper.
%%% Use the second variant for submission (for double-blind reviewing).

\documentclass{ecai} 
%\documentclass[doubleblind]{ecai} 

%%%%%%%%%%%%%%%%%%%%%%%%%%%%%%%%%%%%%%%%%%%%%%%%%%%%%%%%%%%%%%%%%%%%%%%%

%%% Load any packages you require here.
\usepackage{latexsym}
\usepackage{amssymb}
\usepackage{amsmath}
\usepackage{amsthm}
\usepackage{booktabs}
%\usepackage{enumitem}
\usepackage{graphicx}
\usepackage{color}
\usepackage{preamble}

\usepackage{versions}
\usepackage[show]{ed}
%\usepackage[backend=biber, style=authoryear]{biblatex}
% \addbibresource{kwarc.bib}
% \addbibresource{extpubs.bib}
% \addbibresource{kwarccrossrefs.bib}
% \addbibresource{kwarcpubs.bib}

%%%%%%%%%%%%%%%%%%%%%%%%%%%%%%%%%%%%%%%%%%%%%%%%%%%%%%%%%%%%%%%%%%%%%%%%

%%% Define any theorem-like environments you require here.

\newtheorem{theorem}{Theorem}
\newtheorem{lemma}[theorem]{Lemma}
\newtheorem{corollary}[theorem]{Corollary}
\newtheorem{proposition}[theorem]{Proposition}
\newtheorem{fact}[theorem]{Fact}
\newtheorem{definition}{Definition}
\newtheorem{example}[theorem]{Example}

%%%%%%%%%%%%%%%%%%%%%%%%%%%%%%%%%%%%%%%%%%%%%%%%%%%%%%%%%%%%%%%%%%%%%%%%

%%%%%%%%%%%%%%%%%%%%%%%%%%%%%%%%%%%%%%%%%%%%%%%%%%%%%%%%%%%%%%%%%%%%%%%%

%%% Define any new commands you require here.

\newcommand{\BibTeX}{B\kern-.05em{\sc i\kern-.025em b}\kern-.08em\TeX}
\newcommand{\matilda}{\bar{\lambda}\mu \tilde{\mu}}

%%%%%%%%%%%%%%%%%%%%%%%%%%%%%%%%%%%%%%%%%%%%%%%%%%%%%%%%%%%%%%%%%%%%%%%%

\excludeversion{scratch}
\excludeversion{opentypes}
\excludeversion{gpt}
\includeversion{ecai}


\title{Representing Arguments via Control Operators}
\begin{document}

%%%%%%%%%%%%%%%%%%%%%%%%%%%%%%%%%%%%%%%%%%%%%%%%%%%%%%%%%%%%%%%%%%%%%%%%
\begin{ecai}
\begin{frontmatter}

%%% Use this command to specify your submission number.
%%% In doubleblind mode, it will be printed on the first page.

\paperid{123} 

%%% Use this command to specify the title of your paper.

\title{Guidelines for Preparing a Paper for the \\
European Conference on Artificial Intelligence}

%%% Use this combinations of commands to specify all authors of your 
%%% paper. Use \fnms{} and \snm{} to indicate everyone's first names 
%%% and surname. This will help the publisher with indexing the 
%%% proceedings. Please use a reasonable approximation in case your 
%%% name does not neatly split into "first names" and "surname".
%%% Specifying your ORCID digital identifier is optional. 
%%% Use the \thanks{} command to indicate one or more corresponding 
%%% authors and their email address(es). If so desired, you can specify
%%% author contributions using the \footnote{} command.

\author[A]{\fnms{First}~\snm{Author}\orcid{....-....-....-....}\thanks{Corresponding Author. Email: somename@university.edu.}\footnote{Equal contribution.}}
\author[B]{\fnms{Second}~\snm{Author}\orcid{....-....-....-....}\footnotemark}
\author[B,C]{\fnms{Third}~\snm{Author}\orcid{....-....-....-....}} 

\address[A]{Short Affiliation of First Author}
\address[B]{Short Affiliation of Second Author and Third Author}
\address[C]{Short Alternate Affiliation of Third Author}

%%% Use this environment to include an abstract of your paper.

\begin{abstract}
  	We present an encoding of central notions from formal argumentation theory (in particular, logic-based argumentation) into a syntactic variant of the lambda-bar-mu-mu-tilde calculus. Arguments are represented as the open terms of the calculus. Starting from a dialectical interpretation of the calculus in which we regard terms and contexts as contraries, we proceed to define the argumentative relations of rebut, undercut and support. In particular, rebut is closely related to the calculus' cut rule, undercut can be viewed as a special case of the term (context) introduction rule (with a "stashed" term (context) on one side of the command syntax) and support can be captured via an encoding of alternative proofs in the style of Autexier and Sacerdoti-Coen (2006). Reductive versions of the same relations are also investigated. Our approach enables us to capture both the full internal deductive structure of arguments as well as their external argumentation relations in a single format. The resulting argument(ation) proof terms admit rich dialectical and computational interpretations in the spirit of the proofs-as-programs correspondence. In particular, the argumentative attack relations of rebut and undercut have interpretations as control structures. Based on these we define dialectical, argumentative and computational natural language rendering semantics in the style of Sacerdoti-Coen (2005). An implementation based on the Fellowship prover (Kirchner 2007) is outlined.
\end{abstract}
\ednote{TODO: Remove bibliography references.}
\end{frontmatter}

%%%%%%%%%%%%%%%%%%%%%%%%%%%%%%%%%%%%%%%%%%%%%%%%%%%%%%%%%%%%%%%%%%%%%%%%


\end{ecai}

\maketitle
\tableofcontents


\section{Introduction}

\section{The $\matilda$ Calculus - Dialectical Variant}
%\subsection{Language, Typing and Dialectical Interpretation}

We recap the $\matilda$ Calculus and give a variant of its syntax.
Expressions are divided into terms, environments and commands.
Terms are read left-to-right while environments are read right-to-left.
Commands are read either left-to-right or right-to-left depending on how they are bound.\ednote{Check back later.}
The terminology on the expression level gestures at an interpretation of $\matilda$-expressions as program terms and continuation contexts.
Terms are expressions with outer holes that wait to be passed to a context.
Contexts are expressions with inner holes that await a term to be passed to them.
Commands assign terms to contexts.
The $\mu$ binder when binding to the right represents a call-with-current-continuation (or one of the current continuations in the classical case).
Symmetrically, when binding to the left $\mu$ corresponds to a call-with-current-expression (or one of the current expressions in the intuitionistic case).

  \begin{grammar}[Grammar of the $\matilda$-calculus][][gram:terms]
    \firstcasesubtil{$v$}
      {x
       \gralt \lambda x:\pro{\sigma}.v
       \gralt E \circ v\
       \gralt \mu \alpha:\refu{\sigma}.c
       %\gralt \textcolor{blue}{[x]}
      }
      {Terms}

    \firstcasesubtil{$E$}
      {\alpha
       \gralt E.\refu{\sigma}:\alpha\lambda
       \gralt v \circ E
       \gralt c.\pro{\sigma}:x \mu
       %\gralt \textcolor{blue}{[\alpha]}
      }
      {Environments}

     
    \firstcasesubtil{$c$}
      {\rangle v :\contr{\sigma}:  E \langle
       %\gralt \rangle \_ :[\pro{\sigma}]\contr{\sigma}:  E \langle
       %\gralt \rangle v :\contr{\sigma}[\refu{\sigma}]:  \_ \langle
      }
      {Commands}
  

    \firstcasesubtil{\pro{\sigma,\delta}}
      {\pro{ a  }
       \gralt \pro{\sigma \rightarrow  \delta  }
       \gralt \pro{ \sigma - \delta  }
      }
      {Proofs}

    \firstcasesubtil{\refu{\sigma},\refu{\delta}}
      {\refu{a}
       \gralt \refu{\sigma\rightarrow \delta}
       \gralt \refu{ \sigma  - \delta}
      }
      {Refutations}

    \firstcasesubtil{\contr{\sigma}, \contr{\delta}}
      {\contr{a}
       \gralt \contr{\sigma \rightarrow \delta}
       \gralt \contr{\sigma - \delta}
      }
      {Contradictions}
  \end{grammar}

Sequents have the form \emph{Antecedent $\vdash$ Stoup $\dashv$ Succedent}.
Antecedent and Succedent are conjunctive contexts of terms and environments respectively while the stoup contains exactly one active formula - a term, context or command - at any given time.
They are read in the same direction as their elements.
One the type-level, we give a dialectical interpretation following \cite{lovas-struc-normal} and \cite{} \ednote{Cite Harper book}.
In \cite{herbelin-2000} commands are typed by a sequent with an empty stoup mirroring the usual cut rule.
This can be read as a contradiction.
Its typing rule is as follows:

\begin{center}
  \begin{minipage}{0.48\linewidth}
    \begin{prooftree}\label{rule: S-ELIM}
      \symcut{v}{E}{\sigma}
    \end{prooftree}
  \end{minipage}
\end{center}

Instead of cutting the \emph{issue} of the contradiction we want to keep track of it and introduce a type of \emph{pointed contradictions} (bold-face greek letters).
A contradiction arises when a proof meets a refutation.
Correspondingly, terms are typed as proofs (denoted by greek letters) and environments as refutations (denoted by greek-letter shaped holes).
Proofs have the following typing rules (the $\vdash$-symbol is a meta-variable ranging over the pointed contradictions):

\begin{center}
  \begin{minipage}{0.48\linewidth}
    \centering
    \begin{prooftree}\label{axiom:proofs}
      \symaxiom{x}{\sigma}
    \end{prooftree}
  \end{minipage}
  \begin{minipage}{0.48\linewidth}
    \centering
    \begin{prooftree}\label{rule:proimplintro}
      \symproimplintro{x}{\sigma}{v}{\delta}
    \end{prooftree}
  \end{minipage}
\end{center}
\begin{center}
  \smallskip
  \begin{minipage}{0.48\linewidth}
    \centering
    \begin{prooftree}\label{rule:support}
      \symprodiffintro{E}{\delta}{v}{\sigma}
    \end{prooftree}
  \end{minipage}
  \begin{minipage}{0.48\linewidth}
    \centering
    \begin{prooftree}\label{rule: R-ELIM}
      \symprointro{\alpha}{\delta}{c}{\bot}
    \end{prooftree}
  \end{minipage}
\end{center}

These are mirrored by the typing rules for refutations:

\begin{center}
  \begin{minipage}{0.48\linewidth}
    \centering
    \begin{prooftree}\label{axiom:refutations}
      \symmoixa{\alpha}{\delta}
    \end{prooftree}
  \end{minipage}
  \begin{minipage}{0.48\linewidth}
    \centering
    \begin{prooftree}\label{rule: refudiffintro}
      \symrefudiffintro{\alpha}{\sigma}{F}{\delta}
    \end{prooftree}
  \end{minipage}
\end{center}
  \smallskip
\begin{center}
  \begin{minipage}{0.48\linewidth}
    \centering
    \begin{prooftree}\label{rule: S-ELIM}
      \symrefuimplintro{v}{\sigma}{E}{\delta}
    \end{prooftree}
  \end{minipage}
  \begin{minipage}{0.48\linewidth}
    \centering
    \begin{prooftree}\label{rule:undercutelim}
      \symrefuintro{x}{\sigma}{c}{\bot}
    \end{prooftree}
  \end{minipage}
\end{center}

The interplay of proofs and refutations can be viewed through the lens of game semantics as a dialogue between two opposing agents: \emph{Prover} and \emph{Spoiler}.
Prover (Spoiler) can construct proofs (refutations) directly from axioms using the respective connectives introduction rules.
Alternatively they can reason by contradiction: the right-binding (left-binding) $\mu$ binder denotes a challenge to derive a contradiction given an offered refutation (proof).\ednote{TODO: As far as I am aware the \emph{exact} Lorentzen game for the $\matilda$-calculus has not been worked out/published yet, which is surprising since it should not be so difficult. Perhaps a footnote here.}
Restricting the use of proof by contradiction yields different fragments of the system:
In the intuitionistic fragment, proofs by contradiction may not be nested (achieved in the usual way by restricting succedents to at most a single formula);
In the dual-intuitionistic fragment, refutations by contradiction may not be nested (antecedents are restricted to at most a single formula)
In the minimal (dual-minimal) fragment, additionally no tautological refutations (proof) may be used (succedent/antecedent additionally have to hold at least one formula).

Finally, the calculus' expressions can be reduced as follows:

\begin{align*}
  (\rightarrow&) & \rangle \lambda x. v_1 :: v_2 \circ E \langle &\alignarrow   &\rangle v_2 :  : \rangle v_1 :: E \langle . x \mu\langle \\
  (-&) & \rangle E_1 \circ v :: E_2.\alpha \lambda \langle &\alignarrow  &\rangle \mu \alpha. \rangle v :: E_2 \langle :: E_1\langle \\
  (\mu<&) & \rangle \mu \alpha.c :: E \langle &\alignarrow &c[\alpha \leftarrow E]\\
  (>\mu&) & \rangle v :: c.x\mu \langle &\alignarrow &c[x \leftarrow v]                                                       
\end{align*}

If $(\mu<)$-reductions are applied before $(>\mu)$-reductions one gets call-by-value evaluation, otherwise call-by-name evaluation.
Dialectically speaking, the order of evaluation corresponds to the \emph{burden of argument}.
Call-by-value (Call-by-name) evaluation corresponds to burden of proof (refutation): a concrete refutation (proof) gets substituted in and Prover (Spoiler) better had the goods!
On the other hand the substituted refutation's (proof's) evaluation is delayed or may even never occur.
Explicit order of evaluation manipulation can thus model if a proof (refutation) gets to hold \emph{by default}.
We make use of this fact in the next section.

\section{Embedding Argumentation}

The given variant of the $\matilda$-calculus foregrounds its dialectical character as an interplay between proofs and refutations.
This makes it an interesting setting for the study of formal argumentation.
But while proofs and refutations hold with necessity, arguments are defeasible.
Hence we need to relax the notion of proof.

\subsection{Arguments}

In defining arguments we follow deductive argumentation \cite{besnard-review-argum} which views arguments as conditional proofs of the form $\Sigma \vdash_{\mathcal{L}} \varphi$ where the \emph{assumptions} $\Sigma$ and the \emph{conclusion} $\varphi$ are formulae of a given base logic ith consequence relation $\vdash_{\mathcal{L}}$.
In sequent-based variant \cite{},\ednote{TODO: sequent-based argumentation citation.} they then become single conclusion sequents.
We depart from this approach in two ways:
Firstly, unlike deductive argumentation, we do not wish to abstract away from the interior structure of the arguments.
Hence our approach is to define \emph{arguments} as ``open'' terms and \emph{counterarguments} as ``open'' environments.
Secondly, we want to avoid overloading object-language-level variables whose main functionality is hypothetical reasoning.
Unlike premises in conditonal proofs, argument assumptions are not merely hypothetical but closer to (defeasible) knowledge items.
Hence their main route of elimination is not discharge but defeat.
For this reason we define meta-variables as used in the representation of incomplete in the theorem proving literature \ednote{Cite literature and specific examples here.}

This introduces two distinctions from the standard deductive argumentation approach: firstly, in our setting arguments only depend on the assumptions which they use.
In particular, irrespective of the antecedent (succedent) it is not possible to attack an argument on an assumption that does not appear free in it.
Secondly, the same sequent represents many different arguments.
Instead our way of identifying arguments will function via defining a normal form.

Our approach to defining meta-variables is similar to Kirchner's \cite{}, only that we restrict meta-variables to commands.
Let $\conj{\mathcal{V}},\cons{\mathcal{E}}$ and $\ind{\bot}$ be countably infinite sets of meta-variables denoted by gray lower-case latin, greek and hebrew letters, respectively (e.g. $\conj{x}$ and $\cons{\alpha}$ and $\aleph$). We extend the grammar as follows:

\begin{grammar}

  % \firstcasesubtil{$\conj{v}$}
  % {
  % }
  % {Term Meta-variables}
  % \firstcasesubtil{$\cons{E}$}
  % {
  % }
  % {Environment Meta-variables}
  \firstcasesubtil{$c$}
  {...
     \gralt \rangle \conj{v} :\contr{\sigma}:  E \langle
     \gralt \rangle v :\contr{\sigma}:  \cons{E} \langle
     \gralt \rangle \conj{v} :\contr{\sigma}:  \cons{E} \langle
     \gralt \ind{\aleph}
  }{}
\end{grammar}
\ednote{This still needs fixes. Need to introduce all the necessary notations and when metavariable (contexts) are omitted from the notation.}
Metavariables are typed like so:

\begin{center}
  \begin{minipage}{0.48\linewidth}
    \centering
    \begin{prooftree}\label{axiom:proofs}
      \metasymaxiom{x}{\sigma}
    \end{prooftree}
  \end{minipage}
  \begin{minipage}{0.48\linewidth}
    \centering
    \begin{prooftree}\label{rule:proimplintro}
      \metasymmoixa{\alpha}{\sigma}
    \end{prooftree}
  \end{minipage}
\end{center}

\begin{center}
  \begin{minipage}{0.48\linewidth}
    \centering
    \begin{prooftree}\label{axiom:proofs}
      \infer0{V;\bot, \ind{\aleph};\pro{\Sigma}\vdash \ind{\aleph}\dashv \refu{\Delta};E}
    \end{prooftree}
  \end{minipage}
\end{center}

We call the language obtained in this way the $ \Alambda$-calculus (speak ``Argument Calculus'').\ednote{Probably too early here.}
The restriction of meta-variables to commands does not reduce expressiveness: via $\eta$-equivalence, we can always $\mu$-bind a variable and wrap it in a command expression:\ednote{write down all the $\eta$ rules earlier instead.}

\[ v \leftrightarrow_{\eta} \mu \alpha.\rangle v::\alpha \langle\]

At the same time it will enable us to fine-tune the evaluation of arguments: they will be fully $(\rightarrow)$ and $(-)$ reduced with only $\mu$-redexes remaining.
Thus it does not matter if we substitute a subargument for an assumption. \ednote{TODO: Arguments are intuitionistic with holes; counterarguments are dual-intuitionistic with holes. Example: Peirce's law: long form with rebut, undercut}

With these ingredients in place we are finally ready to define our notion of argument:

\begin{definition}[Argument]
  Given a type $t$ an \emph{argument for $t$}  is an intuitionistic $\Alambda$ expression of type $t$ with an arbitrary number of meta-variables. An \emph{argument against $t$} (or a counterargument to $t$) is a dual-intuitionistic $\Alambda$ expression of type $\refu{t}$ with an arbitrary number of metavariables. An \emph{argument} is either an argument for or against some type $t'$.
\end{definition}

The argumentation process is modelled through substitutions of arguments for metavariables.
The basic substitution operation is simply chaining arguments via the rules

\begin{center}
  \begin{minipage}{0.48\textwidth}
    \centering
    % \begin{prooftree}
    %   \hypo{V;\pro{\Sigma}\vdash t[\refu{\sigma}:\alpha] \dashv \refu{\Delta};E}
    %   \hypo{V';\pro{\Sigma}\vdash \refu{\sigma}:E \dashv \refu{\Delta};E'}
    %   \infer2{V,V';\pro{\Sigma}\vdash t[E/\alpha] \dashv \refu{\Delta};E',E}
    % \end{prooftree}
    \envsub
  \end{minipage}
  \begin{minipage}{0.48\textwidth}
    \centering
    % \begin{prooftree}
    %   \hypo{V;\pro{\Sigma}\vdash [X:\pro{\sigma}]t \dashv \refu{\Delta};E}
    %   \hypo{V';\pro{\Sigma}'\vdash v:\pro{\sigma} \dashv \refu{\Delta}';E'}
    %   \infer2{V'';\pro{\Sigma}''\vdash [v/X]t \dashv \refu{\Delta}'';E''}
    % \end{prooftree}
    \termsub
  \end{minipage}
\end{center}
\ednote{Third case for contradictions missing. It is just the same thing again. Make this just one rule for all three cases.}
\ednote{Switch to the cleaner presentation in Kirchner 2007. The inference rules are overkill and much harder to read with all the extra irrelevant symbols.}
where $[v/X]$ denotes the following capturing substitution:

% \begin{itemize}
% \item for any open variable in $v$, $x:\sigma$, if $x$ is in the scope of a binder $\lambda y:\sigma$ or $\mu y: \sigma$ in $[v/x]t$ unify $x$ and $y$,
% \item else if there is a free occurence of $y:\sigma$ in $[v/x]t$, unify $x$ and $y$.
% \end{itemize}


$\pro{\Sigma}''$ is defined as $\pro{\Sigma}\cup \pro{\Sigma}'\setminus \{x\in\pro{\Sigma}:x \text{ was unified in }[v/x]t\}$ and $V''$ is defined as $V\cup V' \setminus \{X\}$

\ednote{Have o3 improve this definition.}

\ednote{Example Peirce's law and dual Peirce's law.}
\begin{example}
  Continuing example \ref{fjf}, Peirce's law can be viewed as the following sequence of argument chainings, applied to an initial argument $\mu*.\rangle \conj{v}: \contr{((\sigma\rightarrow \delta) \rightarrow \sigma) \rightarrow \sigma}: * \langle$ :
  \begin{align*}
    &[ \conj{v} \leftarrow_{\aleph} \lambda f.\mu \alpha.\ind{\aleph}] & S\\
    &[\ind{\aleph} \leftarrow_{v,E} \rangle \conj{v} :\contr{(\sigma\rightarrow \delta) \rightarrow \sigma}: \cons{E} \langle] & P\\
    %&[\conj{v}  \leftarrow_{v'} \mu*.\rangle \conj{v'}:\contr{(\sigma\rightarrow \delta)\rightarrow \sigma}:* \langle] & S\\
    %&[\conj{v'}  \leftarrow f:\pro{(\sigma\rightarrow \delta) \rightarrow \sigma}] & P\\
    &[\conj{v}  \leftarrow f:\pro{(\sigma\rightarrow \delta) \rightarrow \sigma}] & P\\
    &[\cons{E}  \leftarrow_{\ind{\beth}}\lambda x:\sigma.\mu:\delta.\ind{\beth} \circ \alpha] & S\\
    &[\ind{\beth} \leftarrow_{F}\rangle x:\contr{\sigma}:\cons{F}\langle] & P\\
    &[\cons{F} \leftarrow \alpha] & S
  \end{align*}

  In fact, the substitutions can be viewed as a Lorentzen game with the following moves:

  \begin{itemize}
    \itemR{Initial move if $\varphi$ is a proof: $\mu*.\langle \conj{v} :\contr{\varphi}: *\rangle$}{``I claim $\varphi$''}
    \itemR{ Make use of skeptic offer and prove lemma $[\conj{v}  \leftarrow_{V}] $ where $E_1 \circ v_1 \circ \dots \circ v_n$ where $E, v_1 \dots v_n:$  must be either sceptic offers (variables and $\lambda$-captured), $\mu$-captured (challenges) or offers to skeptic ($\refu{\sigma -\delta}$) followed by a challenge ($.\alpha\mu$-wrapped meta-variable $V$).}{``I prove lemma $E_1 \circ v_1 \circ \dots \circ v_n$ using...''.}
  \itemR{Choose a lemma to continue $[\ind{\aleph} \leftarrow_{v,E} \rangle \conj{v} :\contr{(\sigma\rightarrow \delta) \rightarrow \sigma}: \cons{E} \langle]$ where the consequent (leftmost conjunct) must be $\mu$ captured (a challenge).}{``I will prove lemma ... first.''}
  \end{itemize}

Skeptic has exactly the mirrored moves. Hence the dialogue for Peirce's law is rendered as:

\begin{algorithmic}
  \pclaim{$\contr{((\sigma\rightarrow \delta) \rightarrow \sigma)\rightarrow \sigma}$}
     \schall[$\sigma\rightarrow \delta) \rightarrow \sigma$]{$\alpha:\sigma$}
        \pslem{$\contr{((\sigma\rightarrow \delta) \rightarrow \sigma)}$}
           \If {S: I pass.}
              \State \textbf{P:} By $f$.
              \srlem[$x:\sigma$]{$\sigma \rightarrow \delta$}{$\alpha:\sigma$}{$\beta:\delta$}
                  \pslem{$\sigma$}
                     \If {S: I pass.}
                        \State \textbf{P:} By $x$.
                        \srlem{$\top$}{$\alpha:\sigma$}{$\top$}
                        \srmel
                     \Else
                        \srlem{$\top$}{$\alpha:\sigma$}{$\top$}
                        \srmel
                     \EndIf
                  \psmel
              \srmel
           \Else
             \srlem[$x:\sigma$]{$\sigma \rightarrow \delta$}{$\alpha:\sigma$}{$\beta:\delta$}
                  \pslem{$\sigma$}
                     \If {S: I pass.}
                        \State \textbf{P:} By $x$.
                        \srlem{$\top$}{$\alpha:\sigma$}{$\top$}
                        \srmel
                     \Else
                        \srlem{$\top$}{$\alpha:\sigma$}{$\top$}
                        \srmel
                     \EndIf
                  \psmel
              \srmel
           \EndIf
        \psmel
     \sdone
  \pdone
\end{algorithmic}

% P: I claim $\contr{((\sigma\rightarrow \delta) \rightarrow \sigma)}$


\end{example}
\ednote{Need rule that Skeptic cannot break down an offer, else cycles. Add dual Peirce law to the example. Prover home game/away game. Unify naming Skeptic/Spoiler. Should be something stronger than Skeptic since she has her own ambition to refute something. Perhaps R``Refuter''. }
\ednote{Define the long-form that corresponds to these dialogues. Every variable use is mu-abstracted.}
\ednote{I really need to write a Latex exporter for my prover.}


\subsection{Evaluating arguments.}

We extend the $\matilda$ reduction rules to arguments. that capture the intuition of argumentative onus:

\begin{align*}
  (Default <&) & \rangle \mu \alpha.c :: \cons{E} \langle &\alignarrow &c[\alpha \leftarrow E]\\
  (>Default&) & \rangle \conj{v} :: c.x\mu \langle &\alignarrow &c[x \leftarrow v]                                                       
\end{align*}

Intuitively, metavariables behave like opaque terms or \emph{defaults}: they are assumed to represent correctly typed terms and never get evaluated.
\ednote{Coloring metavariables probably left to future work.}



\subsection{Rebut and Cut}

Sequent-based argumentation models argument relations, that is \emph{attacks}, via sequent-elimination rules.
Note the structural similarity between the (compact) rebuttal sequent elimination rule and the $\matilda$ cut rule:

\begin{center}
    \begin{minipage}{0.48\linewidth}
      \centering
      \begin{prooftree}
        \hypo{\Sigma_1\vdash \varphi}
        \hypo{\Sigma_2 \vdash \neg \varphi}
        \infer2{\Sigma_2 \cancel{\vdash} \neg\varphi}
      \end{prooftree}
    \end{minipage}
    \begin{minipage}{0.48\linewidth}
      \centering
      \begin{prooftree}\label{rule: S-ELIM}
      \symcut{v}{E}{\varphi}
    \end{prooftree}
    \end{minipage}
  \end{center}

In the classical version of the $\matilda$-calculus, the rules

\begin{center}
  \begin{minipage}{0.48\linewidth}
    \centering
    \begin{prooftree}
      \hypo{\Sigma \vdash \refu{\varphi}:E \dashv \Delta}
      \infer1{\Sigma \vdash \text{not }E:\pro{\neg\varphi} \dashv \Delta}
    \end{prooftree}
  \end{minipage}
  \begin{minipage}{0.48\linewidth}
    \centering
    \begin{prooftree}
      \hypo{\Sigma \vdash v:\pro{\varphi} \dashv \Delta}
      \infer1{\Sigma \vdash \refu{\neg\varphi}:v \text{ not} \dashv \Delta}
    \end{prooftree}
  \end{minipage}
\end{center}

and their inverses are derivable.
Hence cut and rebut have (semantically) identical application conditions in a classical base logic.
We seek to eliminate this duplication of structure.
However, the rebut rule does not introduce a contradiction but instead eliminates it and introduces an eliminated sequent in one go.
Our task is to encode such sequent elimination using control structures.

\subsection{Encoding Undercut}
\ednote{TODO: starting with reductive versions. If there is still space, add explorative versions.}
Before we tackle the rebut pattern, we will first turn out attention to another mode of attack between arguments: the undercut \ednote{TODO: note on terminology disambiguating defeat, undermine, undercut... Problem: undercut is undermine in aspic, undercut is Pollock's undercutting defeater, defeat itself is not an attack mode but an argument status, and so on.}

  \begin{center}
    \begin{prooftree}
      \hypo{\Sigma_1\vdash \neg \varphi}
      \hypo{\Sigma_2, \varphi \vdash \psi}
      \infer2{\Sigma_2\cancel{\vdash} \psi}
    \end{prooftree}
    \end{center}

Unlike rebut, undercut opererates on an assumption of the attacked argument.

\[undercut(t_1[\conj{v}],t_2) = t_1[\conj{v} \leftarrow \mu \_: \rangle \conj{v} :: t_2\langle ]\]
\ednote{Ednote: Introduce argument notation}

Computationally, undercut throws away any term (environment) that is passed to it (to which it is passed) and instead passes the counterargument $t_2$ to a dedicated error environment (term)  $*$.
In this way the counterargument ``bubbles up'' akin to an exception in a language like OCaml.\ednote{TODO: How can it be caught? Either on the way up encounters catch (on name; on type would require modification of rewrite rule.), i.e. it becomes $\mu{*}.\rangle * :: t_2\langle$; or rebut (a support for the attacked assumption will instead become a rebut). }

\ednote{TODO: undercut argument is labelled OUT.}

\subsection{Encoding Rebut}

We will do so making use of \emph{affine} and \emph{semi-affine} variables.
A $\mu$-bound variable is affine if it does not appear in the command in which it is bound.
It is semi-affine if it only appears on one side of of the command in which it is bound.
Affine variables are either terms that are thrown away or environments that are never passed to.
Exploiting this, affine uses of control operators can direct a term (environment) to a different bound variable.
Semi-affine variables on the other hand are only used in either the call-by-name or call-by-value discipline of execution.
This makes it possible to redirect a term (environment) to a different bound variable conditional on the evaluation order.
In particular, here we define a sentinel variable $\star$ which denotes an error term (environment):

% \begin{align*}
%   rebut_{left}&(t_1,t_2)= \\
%   &\mu *:\refu{\sigma}.\rangle \mu \_ :\refu{\sigma}.\rangle t_1: \contr{\sigma} : * \langle : \contr{\sigma} : \sigma.\rangle\star:\contr{\sigma} t_2:\langle .\pro{\sigma}:\star\mu\langle
% \end{align*}

%\[rebut_{left}(t_1,t_2)= \mu *:\refu{\sigma}.\rangle \mu \_ :\refu{\sigma}.\rangle t_1: \contr{\sigma} :* \langle : \contr{\sigma} : \rangle\star:\contr{\sigma} : t_2\langle .\pro{\sigma}:\_\mu\langle\]
\[rebut_{left}(t_1,t_2)= \mu *:\refu{\sigma}.\rangle \mu \_ :\refu{\sigma}.\rangle t_1: \contr{\sigma} :* \langle : \contr{\sigma} : \rangle\star:\contr{\sigma} : t_2\langle .\pro{\sigma}:\_\mu\langle\]


The reserved environment variable $*$ indicates that the rest of the term is to be evaluated in call-by-value discipline.
Then the term is evaluated in one step as follows:

\begin{align}
  \label{eq:1}
  &\mu *:\refu{\sigma}.\rangle t_1: \contr{\sigma} :* \langle & & (\mu<)
\end{align}


Under normal conditions, this will ensure that the defeated environment $t_2$ is never passed to and instead $t_1$ is eventually returned to any outer context that is substituted for $*$.
But should the burden of proof shift, the term will be reduced as follows:

\begin{align}
  \label{eq:1}
  &\mu *:\refu{\sigma}.\rangle\star:\contr{\sigma} : t_2\langle & & (>\mu)
\end{align}

Now $*$ becomes fully affine and thus the entire outer context gets thrown away if the expression is ever evaluated. Instead the sentinel value $\star$ gets passd to $t_2$.
The proof can now only be completed if a new proof of $\sigma$ is supplied.

In other words, (left) rebut turns a proof into an \emph{argument} predicated on the default assumption, that the counterargument can be defeated in some way.
The burden of refutation lies with Spoiler.
What moves can she make to shift it back to Prover?


\ednote{TODO: rebut without outer $\mu$ binder, two free variables? Can be either captured on the right or on the left. Outer affine variable => undec labelling. How can rebut be relabeled? IN side gets undercut. Then? Or rewriting can only happen if left side not labelled OUT.}


\subsection{Support}

Table captions should be centred \emph{above} the table, while figure 
captions should be centred \emph{below} the figure.\footnote{Footnotes
should be placed \emph{after} punctuation marks (such as full stops).}
 
\begin{table}[h]
\caption{Locations of selected conference editions.}
\centering
\begin{tabular}{ll@{\hspace{8mm}}ll} 
\toprule
AISB-1980 & Amsterdam & ECAI-1990 & Stockholm \\
ECAI-2000 & Berlin & ECAI-2010 & Lisbon \\
ECAI-2020 & \multicolumn{3}{l}{Santiago de Compostela (online)} \\
\bottomrule
\end{tabular}
\end{table}
\subsection{Rebut, Undercut, Support': Reductive Versions}
\section{Argumentation Relations as Control Structures}
\subsection{Recap: The Computational Interpretation of the $\matilda$ Calculus}
\subsection{Rebut, Critical Pairs and Non-determinism}
\subsection{Undercut, Delayed Evaluation and Sentinels}
\subsection{Support and Term Relevance}
\ednote{Support makes sense in domains where terms are mostly irrelevant.}
\section{Explanations in Natural Language}
\section{Related Work}

\begin{scratch}
  \begin{itemize}
  \item Encoding of exceptions in $\lambda\mu$ calculus in $\lambda_{try}$ paper.
  \end{itemize}
\end{scratch}

\section{Conclusion}
Conclusions
\ednote{Prover here}

\begin{scratch}

  \begin{abstract}

    length 100-300 words
    \begin{itemize}
    \item Recently interest in proof theoretic analysis of argumentation
    \item Arguments as computational objects have received less attention
    \item Interactive theorem prover community has made efforts to develop suitable proof languages
    \item This includes representing partial proofs.
    \item Leverage this approach to the representation of arguments.
    \item Take steps towards a proofs-as-programs understanding of argument.
    \item control flow and dialectical interpretation of lambda bar mu tilde mu.
    \end{itemize}
  \end{abstract}
  
  This paper should be about 10 pages, including
  \begin{itemize}
  \item syntax of the calculus
  \item encoding of support, undercut and rebut
  \item possibly: stacked version vs focussed version
  \item Natural language renderings (dialectical, argumentative)
  \end{itemize}

  state of the art:

  \begin{itemize}
  \item what is an argument in ASPIC+?
  \item what is an argument in Carneades?
  \item what is an argument in ABA?
  \item What is an argument in LBA/Sequent-based argumentation?
  \item Arguments in defeasible justification logic
  \item Catala programs?
  \item Declarative programs?
  \item what are natural language representations of arguments in these languages?
  \item what is an argument in proof-theoretic semantics
  \end{itemize}

  Abstractions do not allow to model how the argumentation actually proceeds. Keep history/structure for modelling and explanations. Normalization for

  Computational interpretation of focussed (reductive) undercut: free variable is grabbed by another continuation, its value voided/thrown away and a dummy is passed to the outer continuation, rendering the program worthless.

  future work:

  \begin{itemize}
  \item Rewriting
  \item Labelling/Coloring
  \item Semantics
  \item Default assumptions (Enthymemes)
  \item open types?
  \end{itemize}
  
\end{scratch}

\begin{opentypes}
  \textbf{Notation:} We write $\pro{\varphi}$ for $[\top]\varphi[\bot]$ and correspondingly for $\refu{\varphi},\contr{\varphi}$.
  \begin{grammar}[Adding Open Types to the Grammar][][gram:terms]
    \firstcasesubtil{$v$}
      {x
       \gralt \lambda x:\pro{\sigma}.v
       \gralt E \circ v\
       \gralt \mu \alpha:\refu{\sigma}.c
       %\gralt \textcolor{blue}{[x]}
      }
      {Terms}

    \firstcasesubtil{$E$}
      {\alpha
       \gralt E.\refu{\sigma}:\alpha\lambda
       \gralt v \circ E
       \gralt c.\pro{\sigma}:x \mu
       %\gralt \textcolor{blue}{[\alpha]}
      }
      {Environments}

     
    \firstcasesubtil{$c$}
      {\rangle v :\contr{\sigma}:  E \langle
       %\gralt \rangle \_ :[\pro{\sigma}]\contr{\sigma}:  E \langle
       %\gralt \rangle v :\contr{\sigma}[\refu{\sigma}]:  \_ \langle
      }
      {Commands}
 
    \firstcasesubtil{$h_l$}
      {[\pro{\sigma}]
       %\gralt [\refu{\varphi}]
       \gralt h_l;h_l'
       \gralt [h_l]
      }
      {Left Holes}

    \firstcasesubtil{$h_r$}
      {[\refu{\varphi}]
       \gralt h_r;h_r'
       \gralt [h_r]
       %\gralt [\text{ };\text{ }]
       %\gralt [[\text{ }]]
      }
      {Right Holes} 

    \firstcasesubtil{\pro{\sigma,\delta}}
      {\pro{h_l a h_r }
       \gralt \pro{\sigma \rightarrow h_l \delta h_r }
       \gralt \pro{h_l \sigma - \delta h_r }
      }
      {Programms}

    \firstcasesubtil{\refu{\sigma},\refu{\delta}}
      {h_l\refu{a}h_r
       \gralt h_l\refu{\sigma\rightarrow \delta}h_r
       \gralt h_l\refu{ \sigma  - \delta}h_r
      }
      {Continuations}

    \firstcasesubtil{\contr{\sigma}, \contr{\delta}}
      {h_l\contr{a}h_r
       \gralt h_l\contr{\sigma \rightarrow \delta}h_r
       \gralt h_l\contr{\sigma - \delta}h_r
      }
      {Halt}
  \end{grammar}

\end{opentypes}

\begin{gpt}
  
Below is a compact “fill-in’’ that completes the unfinished paragraph and gives the minimal set-up needed for the rest of the section.  Adapt/expand as you see fit.

----------------------------------------------------------------
\section{The $\matilda$ Calculus – Dialectical Variant}\label{sec:matilda}

We briefly recall the syntax of $\matilda$ (the $\bar\lambda\mu\tilde\mu$-calculus) and sketch the dialectical reading that will ground our argumentation-theoretic encoding.

\subsection{Syntax}

\begin{align*}
  t,u \;::=&\; x 
        \mid \lambda x.t
        \mid \mu \alpha.c
        \mid t\,u\\[2mm]
  e,f \;::=&\; \alpha
        \mid \tilde{\lambda}\alpha.e
        \mid \tilde{\mu}x.c
        \mid t\;::\;e\\[2mm]
  c   \;::=&\; \langle t \mid e\rangle
\end{align*}

Terms $t$ are written left-to-right; environments $e$ are written right-to-left; commands $c$ pair a term with a (co-)context.  Free variables in $t$ are \emph{outer} holes, waiting to be plugged into an environment; free co-variables in $e$ are \emph{inner} holes, waiting for a term.

\paragraph{Binders and Control Operators.}
The binders $\mu$ and $\tilde\mu$ act as control operators:

• $\mu\alpha.c$ captures the \emph{current continuation} (the surrounding environment) and binds it to $\alpha$ inside $c$.  
  Operationally, this is the classical \textsf{call/cc}.  

• Symmetrically, $\tilde\mu x.c$ captures the current \emph{co-continuation} (the surrounding term) and binds it to $x$.

Hence, when we say
\begin{quote}
“The $\mu$-binder when binding to the right represents a \textsf{call/cc}\,,”
\end{quote}
we mean that the evaluation of $\mu\alpha.c$ first stores the ambient context in $\alpha$ and then proceeds with $c$, enabling non-local jumps once $\alpha$ is invoked.

\subsection{Typing}

Types are formulae of the implicational fragment of classical logic:

\[
  A,B ::= p \mid A \to B \mid \bot .
\]

Judgements come in three flavours:

\[ 
  \Gamma \vdash t : A \qquad
  e : A \vdash \Delta \qquad
  \Gamma \vdash c \vdash \Delta .
\]

The rules for $\lambda$, $\mu$ and $\tilde\mu$ mirror the sequent calculus ($\lambda$ for $\to$I, application for $\to$E, $\mu$ for \textsf{cut}-right, $\tilde\mu$ for \textsf{cut}-left).  Full details can be found in \cite{curien2011sym}.

\subsection{Dialectical Reading}

We adopt the bilateralist perspective of Lorenzen games:

• A term $t$ is a \emph{thesis} (a move by Proponent).  
• An environment $e$ is an \emph{antithesis} (a move by Opponent).  
• A command $\langle t \mid e\rangle$ is an interaction (dialogue) in which the thesis meets the antithesis.

Under this reading

– $\beta$-reduction models \emph{persuasion} (rebuttal),  
– $\mu$ / $\tilde\mu$-reductions model \emph{undercut} moves, and  
– alternative normalisation paths correspond to \emph{support} in the sense of \cite{autexier2006altern}.

These correspondences will be made precise in Sect.~\ref{sec:argrels}, after we have defined the three key argument-attack relations in terms of $\matilda$ reduction steps.

----------------------------------------------------------------

This should close the dangling sentence, introduce the syntax, the operational reading of $\mu$, and prepare the ground for the upcoming sections on argument representation.

Below is a compact, self-contained continuation that finishes the incomplete sentence, provides a short “Introduction” and a minimal “Language, Typing, and Dialectical Interpretation” subsection.  Paste it after your current \section{Introduction} and \section{The $\matilda$ Calculus – Dialectical Variant} headers, then delete the placeholders you no longer need.

```latex
%%%%%%%%%%%%%%%%%%%%  Introduction  %%%%%%%%%%%%%%%%%%%%

\section{Introduction}
Argumentation theory offers high–level abstractions for representing and
evaluating conflicting pieces of information, while proof theory offers
low–level, fine-grained representations of logical derivations.  Bridging
these two descriptions promises richer computational interpretations and
a uniform treatment of (i) the *internal* logical structure of arguments
and (ii) their *external* attack/support relations.

λ̄μμ̃ ($\matilda$) is a sequent-style extension of the $\lambda$-calculus
equipped with control operators $\mu$ and $\tilde\mu$.  It supports an
elegant proofs-as-programs reading in which terms correspond to programs
and environments to evaluation contexts.  In this paper we:

\begin{itemize}
  \item encode logic-based arguments as open $\matilda$ terms;
  \item show how \emph{rebut}, \emph{undercut}, and \emph{support} emerge
        naturally from the cut, $\mu$-, and proof-enumeration rules;
  \item derive computational readings in which attack corresponds to
        first-class control operators;
  \item sketch an implementation on top of the Fellowship prover.
\end{itemize}

The remainder is structured as follows.  Section \ref{sec:matilda}
presents a dialectical variant of $\matilda$.  Section 3 formalises the
argumentative relations. Section 4 discusses reductions and semantics,
and Section 5 outlines the implementation.

%%%%%%%%%%%%%%%%%%%%  λ̄μμ̃ Calculus  %%%%%%%%%%%%%%%%%%%%

\section{The $\matilda$ Calculus – Dialectical Variant}\label{sec:matilda}

\subsection{Language, Typing, and Dialectical Interpretation}
\paragraph{Syntax.}
We fix two disjoint, countably infinite sets of variables:
term variables $x,y,z,\dots$ and continuation variables
$\alpha,\beta,\gamma,\dots$.  The grammar is:

\medskip
\begin{minipage}{0.48\linewidth}
\begin{align*}
t &::= x \mid \lambda x.\,t \mid \mu\alpha.\,c \mid t\,t \\
e &::= \alpha \mid e\;t \mid \tilde\mu x.\,c       \\
c &::= \langle t \mid e \rangle
\end{align*}
\end{minipage}
\hfill
\begin{minipage}{0.48\linewidth}
\textit{(terms)}\\
\textit{(environments)}\\
\textit{(commands)}
\end{minipage}
\medskip

Intuitively, $t$ is a program with an \emph{outer} hole waiting for a
continuation, $e$ is a continuation with an \emph{inner} hole expecting
a program, and $c$ wires the two together.  The binder
\smash{$\mu\alpha.\,c$} captures the current continuation—i.e.\ implements
\textsf{call/cc}—while \smash{$\tilde\mu x.\,c$} dually captures the
current program prefix.

\paragraph{Typing.}
Types are standard:

\[
A,B ::= p \mid A \to B,
\]
where $p$ ranges over atomic propositions.  Sequents are of the form
$\Gamma \vdash t : A \mid \Delta$ and $\Gamma \mid e : A \vdash \Delta$
for terms and environments, respectively, with $\Gamma,\Delta$ (multi)-sets
of formulae.  Typing rules mirror those of the classical
λμ-calculus; we only display the distinctive ones:

\[
\frac{\Gamma \vdash t : A \mid \alpha:A,\Delta}{\Gamma \vdash
  \mu\alpha.\,\langle t \mid \alpha\rangle : A \mid \Delta}
\quad
\frac{\Gamma,x:A \mid e : B \vdash \Delta}{\Gamma \mid
  \tilde\mu x.\,\langle x \mid e\rangle : A\to B \vdash \Delta}
\]

\paragraph{Reduction.}
Commands evaluate by
\[
\langle \lambda x.\,t \mid e \rangle \;\longrightarrow\; \langle t[x:=x'] \mid e \rangle
\qquad
\langle t \mid \alpha\rangle \;\longrightarrow\; t[\alpha:=e']
\]
with the usual hygiene side-conditions; the control rules for $\mu$ and
$\tilde\mu$ reflect classical cut elimination.

\paragraph{Dialectical Reading.}
We read $\langle t \mid e\rangle$ as a \emph{dialectical clash} between
contraries $t$ and $e$.  A proof (resp.\ refutation) is a normal form
program (resp.\ continuation).  This symmetry lets us express the
argumentative moves:

\begin{itemize}
\item \emph{Rebut} $A\mapsto\neg A$ is the cut of dual conclusions.
\item \emph{Undercut} attacks a rule by capturing its continuation via
      $\mu$ / $\tilde\mu$.
\item \emph{Support} is modelled by the existence of an alternative
      normalising derivation for the same head conclusion.
\end{itemize}

The next section formalises these ideas.
```

You can, of course, expand or shrink each paragraph depending on space
constraints.
When binding to the left, it corresponds to a call-with-current-value.
\end{gpt}
%%%%%%%%%%%%%%%%%%%%%%%%%%%%%%%%%%%%%%%%%%%%%%%%%%%%%%%%%%%%%%%%%%%%%%%%

%%% Use this command to include your bibliography file.

\bibliography{kwarc,extpubs,kwarccrossrefs,kwarcpubs}
%\printbibliography
\end{document}

% Local Variables:
% gptel-model: o3
% gptel--backend-name: "ChatGPT"
% gptel--bounds: nil
% TeX-master: t
% End:
